% Post-revert state (2026-08-21): only the PART-TWO refactors kept as AI-worded
% ("already written down" topic sentence; attribute-layer insertion + "attribute
% groups"; "taste groups" short form). Moves 1/2/4 reverted to Matteo's own text.
% LLM-usage-disclosure artifact.

An interaction occurs when the characteristics of an item match the taste of a
user, and matrix factorization models mirror that shape when they compute a
recommendation score as the product of a user embedding and an item embedding.
ProtoMF keeps this bilinear form but routes it through prototypes: users and
items are represented by their similarity to a small set of learned prototype
vectors. Each item and user still owns a learned embedding; the prototypes
just act as the structural anchors of the embedding spaces, and every score is
composed from positions relative to them.

Prototypes and Embeddings as is common in MF models are learned exclusively
from the collaborative signal and often at massive scale. The collaborative
signal records who interacted with what, and nothing about what any of those
things are, so no dimension of a representation learned from it carries a
meaning a human could read off. To leverage its prototypes for
interpretability, ProtoMF introduces post-hoc methods to essentially name them
(Section~\ref{sec:bg-protomf}). As discussed in
Section~\ref{sec:bg-iml-intrinsic}, there are issues connected to post-hoc
readings of opaque objects, and where possible, intrinsic methods are
preferred.

We aim to close that step by grounding the prototypes. Grounding, as used in
this thesis, means constraining the latent factors to a chosen interpretable
vocabulary: the model may only build its internal representations from
ingredients that carry meaning on their own. We want to emphasize that the
vocabulary might consist of additional information about elements, but we
still impose it as a constraint for interpretability, not primarily as a
contribution to the model's predictive power. The approach of constraining a
model's internals to expert-defined ingredients (expert vocabulary) has
precedent outside of recommendation, where it was applied likewise for
interpretability at no cost to accuracy. (cite sleep paper, maybe find other
examples.)

A prototype grounded this way explains itself. Its meaning is its composition,
not a name assigned afterwards, and the property propagates to the score. A
recommendation is the user's taste by similarity to Taste-groups, matched with
an item's characteristics by similarity to Item-Characteristics-Groups.

Grounding, however, needs its vocabulary to come from somewhere. What the two
ingredients offer differs sharply, and this is addressed throughout the
lineage of models in this thesis.

For the item ingredient, the vocabulary is already written down. A user
approaches an item with no prior external knowledge; to them, an item simply
is the characteristics shown about it. The catalogue records a coarse subset
of these characteristics as attributes, and to the extent that these
attributes enter the model as interpretable features, they provide the
appropriate vocabulary to explain item prototypes as attribute groups.

Explaining user prototypes as taste groups is not as straightforward. Users do
not state their taste, and in most cases would not be able to articulate it if
asked. Instead, users reveal their taste through their interactions. A CF user
embedding captures this, but as stated, as an unlabeled, uninterpretable taste
vector. User attributes correlate with taste only loosely, are thin in
practice, and connect to privacy concerns when used for explanations; they
therefore enter only as a secondary mechanism (see
Section~\ref{sec:fu-cold-users}). Interaction histories must be the primary
signal, adapted into an interpretable vocabulary.

The next section therefore constructs the easier item side. The harder user
side follows in Chapter~\ref{ch:fu}, where its grounding routes through the
item vocabulary established here, and Chapter~\ref{ch:fufi} merges the two.